\section{Related Work}
\label{sec:related}

\paragraph{KG-augmented LLM reasoning.}
The integration of knowledge graphs with LLMs has matured into a rich research area~\cite{pan2023unifying,han2024graphrag}. Think-on-Graph (ToG)~\cite{sun2023thinkongraph} pioneered LLM-guided interactive exploration of KG neighborhoods, iteratively selecting relations and pruning entities; ToG~2.0~\cite{ma2024thinkongraph} added knowledge-guided retrieval. Graph-constrained Reasoning (RoG)~\cite{luo2024graphconstrained} trains a path-selection model to constrain LLM generation, while GNN-RAG~\cite{mavromatis2025gnnrag} uses graph neural networks to retrieve reasoning paths. Paths-over-Graph~\cite{tan2025pathsovergraph} further refines path search strategies. A common limitation is that these methods treat the KG as immutable and trusted; they do not refine the KG, nor do they verify that final outputs are entailed by specific paths. Recent work on grounding LLM reasoning under incomplete graph evidence~\cite{li2026grounding} and traceability in medical LLM reasoning~\cite{zhu2026hybrid} addresses related concerns but does not combine KG cleaning with post-hoc faithfulness verification. KG-SCoRE differs by closing a verification loop and by reporting faithfulness@$k$.

\paragraph{Literature-based discovery and hypothesis generation.}
Literature-based discovery (LBD), introduced by Swanson~\cite{swanson1986undiscovered}, seeks hidden cross-disciplinary links by chaining literatures~\cite{hristovski2015constructing,bhasuran2023literature}. Classical LBD relies on co-occurrence statistics and graph traversal over SemMedDB~\cite{kilicoglu2012semmeddb,kilicoglu2020broadcoverage}, producing ranked candidates rather than explanatory hypotheses. The LLM era has revitalized LBD: Qi et al.~\cite{qi2024large} systematically evaluated LLMs as biomedical hypothesis generators, Liu et al.~\cite{liu2024literature} combined literature and data signals, and BioVerge~\cite{yang2025bioverge} benchmarked self-evaluating agents. Tyagin et al.~\cite{tyagin2025biomedical} introduced graph-based context retrieval for explainability, and BioHarness~\cite{xiao2026bioharness} assembles evidence across literature, knowledge bases, and atlases. These works focus on retrieval and generation quality but do not co-refine the underlying KG or quantitatively measure faithfulness. KG-SCoRE unifies LBD's symbolic graph heritage with LLM generation under a single auditable loop.

\paragraph{Knowledge graph embedding and completion.}
KG-completion models learn low-dimensional embeddings to score triple plausibility and predict missing links~\cite{ji2021survey,hogan2021knowledge,rossi2021knowledge}. ComplEx~\cite{trouillon2016complex,trouillon2017knowledge} models entities and relations in complex space, capturing antisymmetric patterns; ConvE~\cite{dettmers2018convolutional} applies 2D convolutions, and R-GCN~\cite{schlichtkrull2018modeling} uses graph convolutions. We adopt ComplEx for its efficiency and expressiveness, but repurpose its scores not for link prediction alone---rather, to \emph{gate} LLM verification of suspect triples, a novel use of KG embeddings as a noise detector.

\paragraph{RAG and faithfulness.}
Retrieval-augmented generation grounds LLM outputs in retrieved passages~\cite{brown2025systematic,neha2025retrievalaugmented}. Recent surveys catalog trust frameworks and iterative refinement for faithfulness. Yet RAG faithfulness is typically measured against retrieved \emph{text}, not against explicit \emph{symbolic paths}. Our faithfulness@$k$ metric verifies entailment against KG paths, providing a stronger, structured grounding guarantee than passage-level overlap.
